\documentclass[12pt,a4paper]{article}

\usepackage[utf8]{inputenc}
\usepackage[T2A]{fontenc}
\usepackage[ukrainian]{babel}

\usepackage[
    left=30mm,
    right=15mm,
    top=20mm,
    bottom=20mm
]{geometry}
\usepackage{setspace}
\onehalfspacing

\usepackage{amsmath}
\usepackage{amssymb}
\usepackage{array}
\usepackage{booktabs}
\usepackage{longtable}
\usepackage{listings}
\usepackage{xcolor}

\lstset{
    inputencoding=utf8,
    extendedchars=true,
    basicstyle=\ttfamily\small,
    breaklines=true,
    frame=single,
    columns=fullflexible,
    keepspaces=true
}

\usepackage{graphicx}
\usepackage{float}
\usepackage[
    hidelinks
]{hyperref}
\setlength{\parindent}{1.25cm}
\setlength{\parskip}{0pt}
\setcounter{tocdepth}{1}

\begin{document}


\begin{titlepage}

\begin{center}

КИЇВСЬКИЙ ПОЛІТЕХНІЧНИЙ ІНСТИТУТ
ІМ. ІГОРЯ СІКОРСЬКОГО

\vspace{2cm}

{\large ЗВІТ}\\
{\large ДО ЛАБОРАТОРНОЇ РОБОТИ №1}

\vspace{1cm}

{\large на тему:}

\vspace{0.5cm}

\textbf{«Вибір та реалізація базових фреймворків та бібліотек»}

\end{center}

\vfill

\begin{flushright}
\begin{minipage}{0.4\textwidth}

Виконали:\\
студентки групи ФІ-62 мн\\
Дегтярьова Марія\\
Тимошенко Марія

\end{minipage}
\end{flushright}

\vfill

\begin{center}
Київ -- 2026
\end{center}

\end{titlepage}

\tableofcontents
\newpage

\section{Дослідження принципів реалізації криптографічних механізмів}

\subsection{Чому при розробці прикладного програмного забезпечення не рекомендується самостійно реалізовувати стандартні криптографічні алгоритми?}

% Відповідь

\subsection{Які основні групи криптографічних примітивів та допоміжних функцій зазвичай надають сучасні криптографічні бібліотеки?}

% Відповідь

\subsection{У чому полягає різниця між криптографічним примітивом та криптографічним протоколом? Навести приклади.}

% Відповідь


\subsection{Що таке криптографічно стійкий генератор псевдовипадкових чисел (CSPRNG) та чому звичайні генератори псевдовипадкових чисел не повинні використовуватися для генерації криптографічних ключів?}

% Відповідь


\subsection{У чому полягає різниця між низькорівневими (low-level) та високорівневими (high-level) криптографічними API? Які ризики може створювати неправильне використання низькорівневого API?}

% Відповідь


\subsection{Що означає реалізація криптографічної операції з постійним часом виконання (constant-time implementation) та від яких атак вона повинна захищати?}

% Відповідь


\subsection{Що таке FIPS 140-3 та в яких випадках наявність відповідної сертифікації криптографічного модуля може бути важливою?}

% Відповідь


\subsection{Чому для реалізації RSA та деяких інших асиметричних криптосистем необхідна арифметика багаторозрядних чисел? Які вимоги, крім швидкодії, висуваються до її реалізації у криптографічному програмному забезпеченні?}

% Відповідь
\section{Порівняння криптографічних бібліотек}

\subsection{Обрані криптографічні бібліотеки}

% Коротко описати cryptography та PyCryptodome

\subsection{Порівняльний аналіз бібліотек}

\begin{longtable}{|p{2cm}|p{6cm}|p{6cm}|}
\hline
\textbf{Критерій} &
\textbf{1} &
\textbf{2} \\
\hline

Мови та платформи &
&
\\
\hline

Симетрич\-не шифрування &
&
\\
\hline

AEAD &
&
\\
\hline

Асимет\-рична криптографія &
&
\\
\hline

Цифровий підпис &
&
\\
\hline

Геш-функції та KDF &
&
\\
\hline

Ініціаліза\-ція генератора ви\-падкових чисел &
&
\\
\hline

Крипто\-графічно стійкі випадкові значення &
&
\\
\hline

Високо\-рівневий API &
&
\\
\hline

Апаратне прискорення &
&
\\
\hline

FIPS-сумісність / сертифікація &
&
\\
\hline

Ліцензія &
&
\\
\hline

Вразли\-вості та історія зламів &
&
\\
\hline

Підтримка спільнотою &
&
\\
\hline

Актуаль\-ність документації &
&
\\
\hline

\end{longtable}

\subsection{Обґрунтування вибору бібліотеки}

% Чому обираємо cryptography
\section{Практичне використання криптографічних бібліотек}
\subsection{Симетричне автентифіковане шифрування}

Встановлення бібліотеки.
\begin{lstlisting}[language=Python]
import os
from cryptography.hazmat.primitives.ciphers.aead import AESGCM
from cryptography.exceptions import InvalidTag
\end{lstlisting}

\subsubsection{Генерація криптографічно стійкого ключа}

\begin{lstlisting}[language=Python]
key = AESGCM.generate_key(bit_length=256)
aesgcm = AESGCM(key)

message = b"Hello, my name is Mariia"
\end{lstlisting}

\subsubsection{Шифрування тестового повідомлення}

\begin{lstlisting}[language=Python]
nonce = os.urandom(12)
ciphertext = aesgcm.encrypt(nonce, message, None)

print("Message:", message)
print("Key:", key.hex())
print("Nonce:", nonce.hex())
print("Ciphertext:", ciphertext.hex())

\end{lstlisting}

\subsubsection{Розшифрування повідомлення}

\begin{lstlisting}[language=Python]
decrypted = aesgcm.decrypt(nonce, ciphertext, None)

print("\nDecrypted message:", decrypted)
\end{lstlisting}

\subsubsection{Перевірка цілісності після зміни шифротексту}

\begin{lstlisting}[language=Python]
modified_ciphertext = bytearray(ciphertext)
modified_ciphertext[0] ^= 1
modified_ciphertext = bytes(modified_ciphertext)

print("\nModified ciphertext:", modified_ciphertext.hex())

try:
    aesgcm.decrypt(nonce, modified_ciphertext, None)
    print("Decryption successful")
except InvalidTag:
    print("Error: authentication tag verification failed")
\end{lstlisting}

\subsubsection{Вивід}
У результаті виконання програми було згенеровано 256-бітний ключ, 12-байтовий nonce та отримано шифротекст тестового повідомлення.

\begin{lstlisting}[language=Python]
Message: b'Hello, my name is Mariia'
Key: fc5febdc3df5cd09ecb790ebea5651e7dd441783d988c1e643e58acc8db8f25d
Nonce: 9f72c8ca0d7e717102b1c728
Ciphertext: 1e256814693b724c81fb72b20cb1ee6bcae3db883a2a5fb36416e6af01e720ed626effe2d343afd2

Decrypted message: b'Hello, my name is Mariia'

Modified ciphertext: 1f256814693b724c81fb72b20cb1ee6bcae3db883a2a5fb36416e6af01e720ed626effe2d343afd2
Error: authentication tag verification failed
\end{lstlisting}

Після зміни одного байта шифротексту його розшифрування завершилося помилкою InvalidTag. Це свідчить про виявлення порушення цілісності зашифрованих даних.

\subsubsection{Призначення ключа, nonce та authentication tag}

У використаному алгоритмі AES-GCM ключ є секретним значенням, за допомогою якого виконується шифрування 
та розшифрування повідомлення. У цьому прикладі використовується ключ довжиною 256 біт.

Nonce --- це значення, яке використовується разом із ключем під час шифрування. 
У нашому випадку nonce має довжину 12 байт (Для AES-GCM рекомендована довжина nonce становить 12 байт (96 біт))
і генерується за допомогою криптографічно стійкого генератора випадкових чисел.
Для AES-GCM важливо, щоб nonce не повторювався при використанні одного й того самого ключа. 
Повторне використання nonce з тим самим ключем є небезпечним, оскільки може призвести до розкриття інформації про відкритий текст та порушення автентифікаційного захисту. 
Тому для кожного нового повідомлення необхідно використовувати новий nonce.

Authentication tag --- це тег, який генерується під час шифрування та використовується для 
перевірки цілісності й автентичності шифротексту. У використаній бібліотеці він формується автоматично.

Після зміни одного байта шифротексту перевірка authentication tag не проходить, тому виникає  \texttt{InvalidTag}. 
Це показує,що AES-GCM забезпечує не тільки конфіденційність даних, а й перевірку їх цілісності.

\subsection{Цифровий підпис}

\subsubsection{Генерація ключової пари}

% Код та результат

\begin{lstlisting}[language=Python]
# Bude cod 
\end{lstlisting}


\subsubsection{Створення цифрового підпису}

% Код та результат

\subsubsection{Перевірка цифрового підпису}

% Код та результат

\subsubsection{Перевірка підпису після зміни повідомлення}

% Код та результат

\section{Вибір криптографічної бібліотеки для різних IT-систем}

\subsection{Сценарій А. Backend-сервіс}

\begin{itemize}
    \item Надійність та перевіреність. Сервер працює з конфеденційними даними - тому це обов'язково.
    \item Підтримка сучасних криптографічних алгоритів. Необіхідність, коли мова йде про надійні алгоритми шифрування, гешування та роботи з ключами.
    \item Якісний генератор випадкових чисел. Криптографічна стійка випадковість необхідна для генерації ключів, salt, nonce.
    \item Захист від типових помилок або атак. Ризик повинен бути зменшений.
    \item Оновлення безпеки. Попри гарну схему, важливо регулярне оновлення, щоб швидко виправити виявлені вразливості.
\end{itemize}

\subsection{Сценарій Б. Мобільний застосунок}

\begin{itemize}
    \item Надійне шиврування локальних даних. Якщо зловмисник отрумає доступ до файлів застосунку, то конфіденційна інформація повинна бути збережена.
    \item Низьке споживання пам'яті. Мобільні пристрої більш обмежені, аніж сервери.
    \item Енергоефективність. Надмірна кількість операцій може швидко зменшувати заряд мобільного пристрою.
    \item Безпечне очищення пам'яті. Після використання секретні дані бажано видалити з пам'яті.
    \item Безпечне зберігання ключів. Ключ не можна просто залишити у файлі або в коді застосунку.
\end{itemize}


\subsection{Сценарій В. Сервіс цифрового підпису}

\begin{itemize}
    \item Мінімальний розмір коду (Small Footprint) та відсутність динамічного виділення пам'яті. Бібліотека повинна займати мінімальний обсяг ROM/Flash і RAM, дозволяти статичне виділення буферів та уникати фрагментації обмеженої динамічної пам'яті (heap).
    \item Підтримка апаратних криптографічних блоків мікроконтролера. Можливість задіяти вбудовані в кристал ESP32 апаратні прискорювачі AES, SHA та RSA/ECC для прискорення обчислень та збереження заряду акумулятора.
    \item Використання алгоритмів на еліптичних кривих (ECC). Застосування Ed25519 або ECDSA замість ресурсомісткого RSA, оскільки еліптична криптографія забезпечує високу стійкість за значно меншої довжини ключів (256 біт замість 2048 біт) і менших витрат оперативної пам'яті.
    \item Підтримка апаратного джерела ентропії (TRNG). Для безпечної генерації ключів та nonce бібліотека повинна отримувати випадкові значення безпосередньо з апаратного генератора шуму ESP32, а не використовувати звичайні некриптографічні PRNG.
\end{itemize}

Обраний засіб: Бібліотека \textit{Mbed TLS}, інтегрована до складу \textit{ESP-IDF}, з увімкненою апаратною підтримкою мікроконтролера ESP32.

\subsection{Сценарій Г. Embedded/IoT-пристрій}

\begin{itemize}
    \item Вимога 1
    \item Вимога 2
    \item Вимога 3
    \item Вимога 4
    \item Вимога 5
\end{itemize}

Пояснення вибраних вимог.

\input{5.tex}


\section*{Висновки}
\addcontentsline{toc}{section}{Висновки}

% Тут будуть висновки


\section*{Перелік використаних джерел}
\addcontentsline{toc}{section}{Перелік використаних джерел}

% Тут будуть джерела
\begin{enumerate}

\item
Шось .
\textit{Щось}.
\\
\url{https://shos.tex}

\end{enumerate}

\end{document}